\clearpage
\VerdierPhysicalPageAnchor{0263}%
\begingroup
\small
\setlength{\parindent}{0pt}%
\setlength{\parskip}{0pt}%
\setlength{\emergencystretch}{2em}%
\newcommand{\VerdierTermRow}[3]{%
  \VerdierSemanticAnchor{#1}%
  \par\noindent\hangindent=1.5em\hangafter=1
  #2。\enspace\textup{#3。}\par\smallskip}

\VerdierTermRow{idxt-p263-r01}{複体の第一標準フィルトレーション}{第 III 章、4.1.2}
\VerdierTermRow{idxt-p263-r02}{複体に関する関手の第一スペクトル系列}{第 III 章、4.4.1、4.4.4}
\VerdierTermRow{idxt-p263-r03}{$G$-圏と $G'$-圏の積}{第 I 章、1.3.11}
\VerdierTermRow{idxt-p263-r04}{擬同型}{第 III 章、1.1.3}
\VerdierTermRow{idxt-p263-r05}{符号則}{第 I 章、1.6.6}
\VerdierTermRow{idxt-p263-r06}{$M$ 型の右（それぞれ左）分解}{第 III 章、1.1.4}
\VerdierTermRow{idxt-p263-r07}{Cartan--Eilenberg 入射分解}{第 III 章、4.6.5}
\VerdierTermRow{idxt-p263-r08}{有限長の Cartan--Eilenberg 入射分解}{第 III 章、4.6.5}
\VerdierTermRow{idxt-p263-r09}{$\Sigma$-拡張（関手 $\int_\varphi$）}{第 I 章、2.2.3}
\VerdierTermRow{idxt-p263-r10}{複体のホモトピー圏への関手の $\Sigma$-拡張}{第 I 章、2.5.10}
\VerdierTermRow{idxt-p263-r11}{多加法関手の $\Sigma$-拡張}{第 I 章、2.4.2}
\VerdierTermRow{idxt-p263-r12}{複体のホモトピー圏への関手の射の $\Sigma$-拡張}{第 I 章、2.5.10}
\VerdierTermRow{idxt-p263-r13}{乗法的系の飽和化}{第 II 章、2.2.5}
\VerdierTermRow{idxt-p263-r14}{射の分裂}{第 II 章、1.2.8}
\VerdierTermRow{idxt-p263-r15}{二つのホモトピーの和}{第 I 章、2.5.4}
\VerdierTermRow{idxt-p263-r16}{右（または左）$\mathcal B$-局所化部分圏}{第 II 章、2.3.5}
\VerdierTermRow{idxt-p263-r17}{アーベル圏の導来圏の有界（それぞれ下に有界、上に有界）対象からなる部分圏}{第 III 章、1.2.4}
\VerdierTermRow{idxt-p263-r18}{三角部分圏}{第 II 章、2.1.5}
\VerdierTermRow{idxt-p263-r19}{飽和三角部分圏}{第 II 章、2.1.6}
\VerdierTermRow{idxt-p263-r20}{連結関手列}{第 III 章、3.3.3}
\VerdierTermRow{idxt-p263-r21}{スペクトル対象に付随するスペクトル系列}{第 II 章、4.3.3}
\VerdierTermRow{idxt-p263-r22}{収束スペクトル系列}{第 II 章、4.4.3}
\VerdierTermRow{idxt-p263-r23}{Cartan--Eilenberg スペクトル系列}{第 III 章、4.6.9}
\VerdierTermRow{idxt-p263-r24}{定常スペクトル系列}{第 II 章、4.4.3}
\VerdierTermRow{idxt-p263-r25}{三角部分圏に付随する乗法的系}{第 II 章、2.1.15}
\endgroup
